\label{sec:methodology}

This research employs Design Science Research (DSR) as its primary methodological framework, a choice grounded in the nature of the research objectives and the disciplinary context of software engineering and information systems research. This section justifies this methodological choice, contrasts DSR with alternative approaches, and maps our contributions to established DSR principles.

\subsection{Why Design Science Research}

Design Science Research is a research paradigm that focuses on the creation and evaluation of IT artifacts intended to solve identified organizational problems \cite{hevner2004design}. Unlike behavioral science research, which seeks to develop and verify theories that explain or predict human or organizational behavior, DSR seeks to extend the boundaries of human and organizational capabilities by creating new and innovative artifacts \cite{hevner2004design, march1995design}.

The fundamental premise of DSR is that knowledge and understanding are achieved through the \textit{building and application} of designed artifacts \cite{simon1996sciences}. This ``knowing through making'' epistemology is particularly suited to software engineering research, where the primary output is often a software system, tool, or methodology rather than a descriptive theory.

\subsection{DSR in Software Engineering and Information Systems}

Design Science Research has emerged as one of the most respected and widely adopted research methodologies in software engineering and information systems fields for several compelling reasons:

\textbf{1. Artifact-Centric Nature of the Discipline}: Software engineering is inherently an applied discipline concerned with the construction of software artifacts. DSR's focus on creating and evaluating artifacts aligns naturally with the discipline's core activities \cite{heistracher2004design}. The Association for Information Systems (AIS) has explicitly recognized DSR as a legitimate and valuable research approach alongside behavioral science \cite{hevner2004design}.

\textbf{2. Problem-Solution Orientation}: DSR addresses wicked problems---complex, ill-structured problems that resist straightforward solutions \cite{rittel1973dilemmas}. In software engineering, many challenges (scalability, maintainability, usability) are inherently wicked, making DSR's iterative design-build-evaluate cycle particularly appropriate.

\textbf{3. Rigorous Evaluation Framework}: Hevner et al. (2004) established seven guidelines for conducting and evaluating DSR, providing a rigorous framework that ensures research quality and relevance \cite{hevner2004design}. These guidelines have become the de facto standard for DSR in information systems.

\textbf{4. Dual Contribution Model}: DSR requires contributions to both theory (design knowledge) and practice (usable artifacts), ensuring that research has tangible impact \cite{venable2016design}. This dual contribution aligns with funding agencies' and industry's expectations for applicable research outcomes.

\textbf{5. Publication Outlets}: Major journals in software engineering and information systems (MIS Quarterly, Information Systems Research, IEEE Transactions on Software Engineering, Journal of Systems and Software) actively publish DSR studies, indicating the field's acceptance and respect for this methodology \cite{straub2015editor}.

\subsection{Comparison with Alternative Methodologies}

To justify DSR as the most appropriate methodology for this research, we compare it with two primary alternatives: Qualitative Research and Quantitative Research.

\subsubsection{Qualitative Research Approaches}

Qualitative research methods---including case studies, ethnography, grounded theory, and phenomenology---focus on understanding phenomena in depth through non-numerical data \cite{myers2019qualitative}.

\textbf{Strengths}: Qualitative approaches excel at exploring complex social phenomena, generating rich descriptions, and developing theoretical insights from empirical observations. They are particularly valuable for understanding \textit{why} and \textit{how} questions.

\textbf{Limitations for This Research}: While qualitative methods could explore how practitioners use contagion simulation tools or what challenges they face, these approaches do not inherently involve \textit{creating} solutions. Our research objective is not merely to understand the problem space but to construct a novel artifact that addresses identified gaps. Purely qualitative inquiry would leave the solution space unexplored.

\subsubsection{Quantitative Research Approaches}

Quantitative research employs statistical and mathematical methods to test hypotheses and establish generalizable findings through numerical data \cite{cohen2018research}.

\textbf{Strengths}: Quantitative methods provide rigorous hypothesis testing, enable generalization through statistical inference, and allow for precise measurement of variables and relationships.

\textbf{Limitations for This Research}: Quantitative approaches typically require pre-existing theories and instruments to test. Our research aims to create a novel artifact---a contagion simulation system integrating cellular automata with agent-based modeling---where no prior theory or instrument exists. Furthermore, quantitative evaluation (e.g., performance benchmarks, user studies) is a \textit{component} of DSR but not the complete methodology.

\subsubsection{Why DSR is Most Appropriate}

DSR subsumes and integrates elements of both qualitative and quantitative approaches while adding the crucial dimension of artifact creation:

\begin{itemize}
    \item \textbf{Problem Understanding (Qualitative Element)}: DSR begins with problem identification and motivation, often employing qualitative techniques to understand stakeholder needs and contextual factors \cite{peffers2007design}.
    
    \item \textbf{Artifact Construction (Design Element)}: The core of DSR is the creation of an innovative artifact---the component absent from pure qualitative and quantitative approaches. This research constructs a spatial contagion sandbox extending the ChatDev multi-agent platform.
    
    \item \textbf{Evaluation (Quantitative Element)}: DSR requires rigorous evaluation, which may include quantitative measures (performance metrics, simulation accuracy), qualitative assessments (expert reviews, user feedback), or mixed methods \cite{hevner2004design}.
    
    \item \textbf{Knowledge Contribution}: DSR generates design knowledge (design principles, patterns, theories) that extends beyond the specific artifact, contributing to the broader body of software engineering knowledge \cite{gregor2007nature}.
\end{itemize}

For research whose primary objective is to \textit{design, implement, and evaluate} a novel software artifact addressing a practical problem, DSR is not merely appropriate---it is the canonical methodological choice \cite{wieringa2014design}.

\subsection{DSR Guidelines and This Research}

Hevner et al. (2004) articulate seven guidelines for effective DSR \cite{hevner2004design}. We map our research to these guidelines to demonstrate methodological rigor:

\textbf{Guideline 1: Design as an Artifact}. DSR must produce a viable artifact in the form of a construct, model, method, or instantiation.

\emph{Our Contribution}: This research produces an instantiation artifact---the spatial contagion sandbox---implemented as a JavaScript-based simulation engine integrated with the ChatDev/DevAll platform. The artifact includes the \texttt{useContagionEngine.js} composable, the PixiJS-based \texttt{SpatialCanvas} environment, and a configuration-driven parameter system.

\textbf{Guideline 2: Problem Relevance}. The objective of DSR is to develop technology-based solutions to important and relevant business problems.

\emph{Our Contribution}: The research addresses the problem of accessible contagion simulation for education and public health training. In the post-COVID-19 era, demand for tools that enable non-programmers to explore epidemiological scenarios has increased dramatically. Existing agent-based epidemiology platforms (GAMA, MASON, NetLogo) require programming expertise, creating barriers for educators and public health professionals. Our artifact lowers these barriers through a zero-code, visual interface.

\textbf{Guideline 3: Design Evaluation}. The utility, quality, and efficacy of a design artifact must be rigorously demonstrated via well-executed evaluation methods.

\emph{Our Contribution}: Evaluation includes: (1) functional correctness verification through scenario testing (hospital and school outbreak simulations), (2) emergent behavior analysis demonstrating the artifact's capacity to produce realistic epidemic dynamics, (3) comparison with established agent-based platforms on dimensions of accessibility and integration capabilities, and (4) informal user feedback from educational deployments.

\textbf{Guideline 4: Research Contributions}. Effective DSR must provide clear and verifiable contributions in the areas of the design artifact, design foundations, and/or design methodologies.

\emph{Our Contribution}: This research contributes: (1) a novel integration of cellular automata (Conway's Game of Life) with agent-based modeling for contagion simulation, (2) a composable architecture pattern for embedding simulation engines within multi-agent platforms, and (3) a configuration-driven approach to epidemiological parameter specification enabling domain experts to customize simulations without code modification.

\textbf{Guideline 5: Research Rigor}. DSR relies upon the application of rigorous methods in both the construction and evaluation of the design artifact.

\emph{Our Contribution}: The artifact construction applies software engineering best practices: modular design through composables, event-driven architecture for real-time rendering, and schema-validated configuration. Evaluation employs established metrics from agent-based modeling literature and follows patterns for computational artifact evaluation \cite{sonnenberg2012evaluation}.

\textbf{Guideline 6: Design as a Search Process}. The search for an effective artifact requires utilizing available means to reach desired ends while satisfying laws in the problem environment.

\emph{Our Contribution}: The design process explored multiple implementation alternatives: (1) cellular automata vs. continuous space models, (2) PixiJS vs. Canvas API for rendering, (3) monolithic vs. composable architecture, and (4) hardcoded vs. configuration-driven parameters. The final artifact represents a satisficing solution balancing computational efficiency, visual fidelity, and extensibility.

\textbf{Guideline 7: Communication of Research}. DSR must be presented effectively both to technology-oriented and management-oriented audiences.

\emph{Our Contribution}: This paper communicates to: (1) software engineering researchers through formal specification of the agent state machine and transmission mechanisms, (2) epidemiology researchers through alignment with established compartmental models (SIR/SEIR), and (3) practitioners through reproducibility artifacts (open-source code, configuration examples, documentation).

\subsection{Design Science Research Process}

Our research follows the Design Science Research Process (DSRP) model proposed by Peffers et al. (2007) \cite{peffers2007design}:

\begin{enumerate}
    \item \textbf{Problem Identification and Motivation}: Identified the gap in accessible, LLM-integrated contagion simulation tools through literature review and analysis of existing platforms (US-002, US-005).
    
    \item \textbf{Define Objectives for a Solution}: Specified objectives: zero-code configuration, real-time visualization, emergent behavior demonstration, and integration with existing multi-agent framework (US-002).
    
    \item \textbf{Design and Development}: Constructed the spatial contagion sandbox artifact, including the state machine, transmission mechanisms, and visual interface (US-008, US-009, US-010).
    
    \item \textbf{Demonstration}: Demonstrated the artifact through hospital and school outbreak scenarios with configured parameters (US-012).
    
    \item \textbf{Evaluation}: Evaluated emergent behaviors, compared with existing tools, and assessed alignment with educational use cases (US-011, US-013).
    
    \item \textbf{Communication}: Communicated results through this research paper, open-source repository, and documentation (all sections).
\end{enumerate}

\subsection{Positioning in DSR Taxonomy}

Following the DSR taxonomy of Gregor and Hevner (2013) \cite{gregor2013positioning}, this research is classified as:

\begin{itemize}
    \item \textbf{Purpose}: Improving an existing capability (contagion simulation) through a novel approach (cellular automata integration)
    \item \textbf{Artifact Type}: Instantiation (working software system)
    \item \textbf{Contribution Type}: Improvement and exaptation (applying cellular automata concepts to agent-based epidemiology in a new way)
    \item \textbf{Evaluation Type}: Artificial, formative (controlled scenario testing rather than field deployment)
\end{itemize}

This positioning acknowledges that the artifact is a working prototype demonstrating feasibility and design principles, with full field deployment and summative evaluation reserved for future work.

\subsection{Limitations and Threats to Validity}

We acknowledge methodological limitations:

\textbf{Construct Validity}: The artifact addresses a specific conceptualization of contagion simulation; alternative conceptualizations (e.g., differential equation-based SIR models) may yield different design decisions.

\textbf{Internal Validity}: Emergent behaviors observed in simulation may reflect implementation artifacts rather than genuine epidemiological dynamics. We mitigate this through comparison with established models.

\textbf{External Validity}: The artifact has been tested with small populations (10-50 agents). Generalization to larger-scale simulations requires further evaluation.

\textbf{Reliability}: As a design artifact, reproducibility depends on exact environment configuration. We address this through comprehensive documentation and containerization (US-007).

\subsection{Summary}

Design Science Research provides the most appropriate methodological framework for this research because: (1) it aligns with the artifact-centric nature of software engineering research, (2) it integrates problem understanding (qualitative) and rigorous evaluation (quantitative) with artifact creation, (3) it provides established guidelines ensuring research quality, and (4) it enables contributions to both theory and practice. By adhering to Hevner et al.'s seven guidelines and following the DSRP process, this research ensures methodological rigor while addressing a relevant, real-world problem through innovative artifact construction.
